Обработка больших данных является одной из ключевых задач современной аналитики. При работе с наборами, содержащими миллионы записей и десятки признаков, последовательная обработка на одном вычислительном узле может требовать значительных временных и аппаратных ресурсов. Для распределённой обработки таких данных применяются Apache Hadoop и Apache Spark, позволяющие хранить данные в распределённой файловой системе и выполнять вычисления параллельно на нескольких узлах кластера.

Актуальность работы обусловлена необходимостью получать достоверные аналитические результаты на больших и неоднородных наборах данных. Данные о дорожно-транспортных происшествиях содержат сведения о времени и месте аварии, погодных условиях, характеристиках дорожной инфраструктуры и тяжести происшествия. При этом исходный набор включает пропуски, аномальные значения и категориальные признаки с высокой кардинальностью. Поэтому перед построением моделей машинного обучения требуется выполнить подготовку и исследование данных, а затем оценить возможности распределённого обучения.

В связи с этим целью данной работы являлось исследование возможностей фреймворка Apache Spark для распределённой обработки и анализа больших данных на примере набора сведений о дорожно-транспортных происшествиях в США, а также оценка применимости методов машинного обучения для решения задач регрессии и классификации.

Для достижения поставленной цели решались следующие задачи:
\begin{enumerate}
    \item изучить особенности больших данных и программные средства их распределённого хранения и обработки;
    \item охарактеризовать набор данных US Accidents и определить состав признаков, используемых в исследовании;
    \item загрузить исходные данные в Apache Spark, выполнить преобразование типов и сохранить подготовленный набор в формате Apache Iceberg;
    \item провести разведочный анализ, выявить пропуски, выбросы и аномальные значения, а также выполнить необходимую обработку признаков;
    \item сформулировать задачи регрессии и бинарной классификации и реализовать конвейеры подготовки признаков и обучения моделей средствами Spark MLlib;
    \item оценить качество построенных моделей на тестовых выборках и сопоставить полученные результаты с базовыми предсказателями.
\end{enumerate}

Объектом исследования является набор данных о дорожно-транспортных происшествиях в США за период с 2016 по 2023 год. Предмет исследования -- методы распределённой обработки, разведочного анализа и машинного обучения, реализованные с использованием Apache Spark.

В первой главе рассматриваются постановка задачи разведочного анализа, структура исходного набора US Accidents, используемое программное окружение, загрузка данных в Apache Spark и основные этапы их очистки и преобразования. Во второй главе рассматриваются задачи регрессии и бинарной классификации, построение конвейеров Spark MLlib, обучение моделей и анализ качества полученных результатов. В заключении сформулированы основные выводы по выполненной работе и обозначены направления дальнейшего исследования.